% Part: proof-theory
% Chapter: natural-deduction
% Section: introduction

\documentclass[../../../include/open-logic-section]{subfiles}

\begin{document}

\olfileid{pt}{nat}{int}

\olsection{تمهيد في الاستنباط الطبيعي}

الاستنباط الطبيعي اسم لعائلة من أنظمة~!!{proof}. ولكل رابط منطقي أو سور
فيها استدلالان متقابلان: قاعدة~!!a{introduction} «تدخل» العامل، وقاعدة
!!a{elimination} «تحذفه». فنتيجة~$\Intro{\lif}$ مثلًا~!!{formula}s من
الصورة~$!A \lif !B$، وفي مقدمات~$\Elim{\lif}$~!!{formula}s من الصورة
$!A \lif !B$.

وأول صورتين للاستنباط الطبيعي وضعهما غيرهارد غنتسن وستانيسواف
ياشكوفسكي في ثلاثينيات القرن العشرين. والذي يغلب في مباحث نظرية البرهان
نظام غنتسن، وأما نظام ياشكوفسكي فمنه نشأت أكثر الأنظمة استعمالًا في
التدريس. ويبدأ كل منهما بـ\emph{فروض}: وهي~!!{formula}s لا يلزم أن
!!{prove}، ولكنها تستعمل في~!!{proving}~!!{formula}s أخرى. وقد يقوم
!!{proof} كله على مثل هذه الفروض؛ ولأنها لم تبرهن، فليس~!!{proof} مثبتًا
لنتيجته، أعني~\emph{!!{formula} نهايته}، وإنما يثبت لزوم النتيجة عن
الفروض.

ومن قواعد الاستدلال ما تنقطع نتيجته عن الاعتماد على بعض الفروض، فيقال
إن القاعدة «تبرئ» تلك الفروض. فـ~$\Intro{\lif}$ تأخذ مثلًا~!!a{proof}
للنتيجة~$!B$ من الفرض~$!A$، وتنتج~$!A \lif !B$. وحين ينتهي~!!{proof}
بـ$!A \lif !B$ لا يبقى معتمدًا على~$!A$، لأن الفرض~$!A$ قد برئ. وفي
أنظمة غنتسن ترسم~!!{proof}s أشجارًا من~!!{formula}s، تكون الفروض فيها
أعلى~!!{formula}s، وتوسم قواعد من قبيل~$\Intro{\lif}$ بما يدل على
الفروض المبرأة. وأما أنظمة ياشكوفسكي فتفصل على وجه ما أقسام~!!{proof}
المعتمدة على فرض: فربما جعلتها وراء خط رأسي أو داخل صندوق؛ ويكون الفرض
$!A$ الذي يراد تبرئته أول أسطر الصندوق، ويكون تالي الشرط~$!B$ آخرها،
وتقع نتيجة~$\Intro{\lif}$، وهي~$!A \lif !B$، خارج الصندوق.

وسميت هذه الأنظمة «طبيعية» لأن قواعدها تجري على طرائق الاستدلال المألوفة،
ولا سيما عند الرياضيين. فإذا أردنا إثبات نتيجة شرطية افترضنا المقدّم
وبرهنا به التالي؛ وهذه~$\Intro{\lif}$. وإذا علمنا الشرط~$!A \lif !B$
وعلمنا مقدمه~$!A$، لزم~$!B$؛ وهذه~$\Elim{\lif}$. وإذا أقمنا البرهان
على الحالات أثبتنا لزوم النتيجة من افتراض الفصل~$!A \lor !B$؛ وهذه
$\Elim{\lor}$. وإذا أردنا دعوى عامة~$\lforall[\OLId{latin-ordinary}{x}][!A(\OLId{latin-ordinary}{x})]$ برهنا~$!A(\OLId{latin-ordinary}{c})$
لثابت اعتباطي~$\OLId{latin-ordinary}{c}$؛ وهذه~$\Intro{\lforall}$. وعلى هذا القياس.

ومن أمثلته~!!{proof} يسير لـ$!A \lif (!A \lor !B)$ في الاستنباط
الطبيعي على رسم غنتسن:
\[
\AxiomC{$\Discharge{!A}{\OLId{latin-ordinary}{x}}$}
\RightLabel{\Intro{\lor}}
\UnaryInfC{$!A \lor !B$}
\DischargeRule{\Intro{\lif}}{\OLId{latin-ordinary}{x}}
\UnaryInfC{$!A \lif (!A \lor !B)$}
\DisplayProof
\]
مبدؤه الفرض~$!A$، ومن~$\Intro{\lor}$ تنتج~$!A \lor !B$، ثم من
$\Intro{\lif}$ تنتج~$!A \lif (!A \lor !B)$. واستدلال~$\Intro{\lif}$
يبرئ الفرض~$!A$، وعلى ذلك تدل العلامة~$\OLId{latin-ordinary}{x}$.

والغالب على منظّري البرهان إيثار أنظمة غنتسن الأصلية المبنية من أشجار
!!{formula}s، وإن كان لأحد تنويعاتها شأن في نظرية البرهان أيضًا. فإن
!!^a{proof} لـ$!A$ من الفروض~$\OLId{greek-variable}{Gamma}$ يدل على لزوم~$!A$ من~$\OLId{greek-variable}{Gamma}$،
أي~$\OLId{greek-variable}{Gamma} \Entails !A$. ولهذا تقرأ قواعد الاستنباط الطبيعي باعتبارها
أحكامًا على علاقات اللزوم؛ فـ~$\Intro{\lif}$ تقول مثلًا إنه إذا كان
$\OLId{greek-variable}{Gamma} + !A \Entails !B$ كان~$\OLId{greek-variable}{Gamma} \Entails !A \lif !B$. ومن ثم
يحسن أن تصاغ القواعد فتتناول، في مقدماتها ونتيجتها، الفروض وما يلزم
عنها جميعًا؛ أي تتناول المتتاليات~$\OLId{greek-variable}{Gamma} \Sequent !A$. ويصير البرهان
اليسير المتقدم في مثل هذا النظام على الصورة الآتية:
\[
\Axiom$\OLId{latin-ordinary}{x}:!A \fCenter !A$
\RightLabel{\Intro{\lor}}
\UnaryInf$\OLId{latin-ordinary}{x}:!A \fCenter !A \lor !B$
\DischargeRule{\Intro{\lif}}{\OLId{latin-ordinary}{x}}
\UnaryInf$\fCenter !A \lif (!A \lor !B)$
\DisplayProof
\]
والقواعد، كما ترى، لم تتغير: فهي تعمل في~!!{formula} الواقعة عن يمين
$\Sequent$، وأما~!!{formula}s التي عن يساره فوظيفتها أن تحصي الفروض
المعتمد عليها في كل خطوة.

لكن أزواج قواعد~!!{introduction} و!!{elimination}، وإن أخذت كلها، لا
تستقل بإتمام المنطق الكلاسيكي؛ فلا بد من قاعدتين تتصلان بـ$\lfalse$.
إحداهما «قاعدة المحال الحدسية»~$\FalseInt$، وتسمى «الانفجار»، ومؤداها
أن كل شيء يستنتج من~$\lfalse$. والثانية مبدأ البرهان الكلاسيكي غير
المباشر~$\FalseCl{}$، أو «قاعدة المحال الكلاسيكية»: فإذا أمكن أن
!!{prove}~$\lfalse$ من~$\lnot !A$ أمكن من ثم أن~!!{prove}~$!A$. فإن
حذفنا~$\FalseCl$ بقي نظام يلائم المنطق الحدسي، وإن حذفنا القاعدتين كان
النظام ما يسمى «المنطق الأدنى».

وسنعالج الاستنباط الطبيعي على طريقة غنتسن في المنطقين الكلاسيكي
والحدسي، أعني النظامين~\Log{N1c} و\Log{N1i}. ويقابلهما~\Log{N2c}
و\Log{N2i}، وهما يعملان في المتتاليات~$\OLId{greek-variable}{Gamma} \Sequent !A$ بدل
!!{formula}s المفردة~$!A$.

\end{document}
